% ===========================================================
%  第 9 章 二次型 —— 高等代数【深化】拓展框（补充片）
%  源：北大《高等代数》第五版 第五章（印刷 138–156）
%      樊启斌《高等代数典型问题与方法》§8.2–§8.3（印刷 511–545）
%  与主控手写块 deep_ch09_authored.tex 互补，不重复其 4 块
% ===========================================================

% ---------- 惯性定理的证明思路 ----------
% 来源：北大 五.3 唯一性（印刷页 149–151）
\begin{fdeep}{惯性定理的证明思路：为什么正负惯性指数是不变量}
同一个实二次型可以化成许多不同的标准形，但规范形只有一个：
\fml{f=d_1y_1^{2}+\dots+d_ry_r^{2}\ (d_i\ne0)\ \Longrightarrow\
f=z_1^{2}+\dots+z_p^{2}-z_{p+1}^{2}-\dots-z_r^{2}}
（实数域中正数总可开平方，故可把每个 $d_i$ 变成 $\pm1$）。关键只在证明 $p$ 唯一。
\textbf{反证法}：设同一 $f$ 经 $X=BY$ 与 $X=CZ$ 分别化成规范形，要证 $p=q$。若 $p>q$，
令 $G=C^{-1}B=(g_{ij})$，则 $Z=GY$。作齐次线性方程组（前 $q$ 个方程取自 $Z=GY$）：
\fml{\begin{cases}
g_{11}y_1+g_{12}y_2+\dots+g_{1n}y_n=0,\\
\dots\dots\dots\dots\dots\dots\\
g_{q1}y_1+g_{q2}y_2+\dots+g_{qn}y_n=0,\\
y_{p+1}=0,\ \dots,\ y_n=0.
\end{cases}}
它含 $n$ 个未知量、而方程只有 $q+(n-p)=n-(p-q)<n$ 个，故有非零解 $(k_1,\dots,k_n)$。
把该解代入规范形的左端，得 $k_1^{2}+\dots+k_p^{2}>0$；代入右端，因它是方程组前 $q$ 个方程的解
故 $z_1=\dots=z_q=0$，得 $-z_{q+1}^{2}-\dots-z_r^{2}\le0$。矛盾。所以 $p\le q$，同理 $q\le p$，
于是 $p=q$，规范形唯一。

【反哺点】：服务考纲〔二次型〕"了解…标准形规范形概念及惯性定理"。"方程个数少于未知量个数 $\Longrightarrow$ 有非零解"是惯性定理的\textbf{唯一考法}：凡是要证"两个由二次型符号定义的数相等"，都从齐次方程组入手；它也顺带说明了正负惯性指数为什么是\textbf{合同不变量}（只依赖二次型本身，与替换无关）。
\end{fdeep}

% ---------- 标准形与规范形的唯一性 ----------
% 来源：北大 五.3（印刷页 148–149）；樊 8.2.1（印刷页 511）
\begin{fdeep}{规范形唯一、标准形不唯一：复域与实域两套规范形}
同一个二次型用不同的非退化线性替换，得到的标准形系数可以完全不同：例如某个三元二次型
既可化为 $2y_1^{2}-2y_2^{2}+6y_3^{2}$，也可化为 $2w_1^{2}-\frac{1}{2}w_2^{2}+\frac{2}{3}w_3^{2}$。
所以\textbf{标准形不唯一}，能唯一确定下来的只是"非零系数的个数"和"正负系数的个数"。
把非零系数继续用 $y_i=\frac{1}{\sqrt{d_i}}z_i$ 化成 $\pm1$，得到两套\textbf{唯一}的规范形：
复二次型的规范形为
\fml{z_1^{2}+z_2^{2}+\dots+z_r^{2}}
实二次型的规范形为
\fml{z_1^{2}+\dots+z_p^{2}-z_{p+1}^{2}-\dots-z_r^{2}}
（依据是复数总能开平方、正实数总能开平方。）由此得到两条\textbf{合同判据}：
复对称矩阵 $A$ 与 $B$ 合同 $\iff$ 秩相等；
实对称矩阵 $A$ 与 $B$ 合同 $\iff$ 正、负惯性指数分别相等（等价于秩相同且正惯性指数相同）。

【反哺点】：服务考纲〔二次型〕"了解二次型秩、合同变换与合同矩阵、标准形规范形概念及惯性定理"。考场上问"两个二次型的标准形是否相同"，答"标准形相同"是\textbf{错}的：标准形不唯一；只有\textbf{规范形}才唯一，而"规范形相同 $\iff$ 秩相同且正惯性指数相同"。这一句把"标准形 / 规范形"这对极易混淆的概念一次性分清。
\end{fdeep}

% ---------- 配方法的两种情形 ----------
% 来源：北大 五.2 定理 1 的证明（印刷页 141–143）；樊 8.2.3(1)（印刷页 512）
\begin{fdeep}{配方法的两种情形与操作步骤}
配方法的本质是把一个变量一次配完、再对剩下的变量重复（对变量个数作归纳）。
\textbf{情形一：有平方项，不妨 $a_{11}\ne0$。} 把含 $x_1$ 的项全部集中配成完全平方：
\fml{f=a_{11}\Bigl(x_1+\sum_{j=2}^{n}a_{11}^{-1}a_{1j}x_j\Bigr)^{2}
+\sum_{i=2}^{n}\sum_{j=2}^{n}b_{ij}x_ix_j}
后半部分只含 $x_2,\dots,x_n$，是 $n-1$ 元二次型。取替换
$y_1=x_1+\sum_{j\ge2}a_{11}^{-1}a_{1j}x_j,\ y_2=x_2,\ \dots,\ y_n=x_n$ 即可。
该替换的系数矩阵是单位上三角阵、行列式为 $1$，故一定\textbf{非退化}。
\textbf{情形二：没有平方项（所有 $a_{ii}=0$）但有交叉项 $a_{ij}\ne0\ (i\ne j)$。}
先"制造"一个平方项，不妨设 $a_{12}\ne0$，取
\fml{x_1=z_1+z_2,\qquad x_2=z_1-z_2,\qquad x_k=z_k\ \ (k=3,\dots,n)}
则 $f=2a_{12}z_1^{2}-2a_{12}z_2^{2}+\dots$，其中 $z_1^{2}$ 的系数 $2a_{12}\ne0$，化为情形一。
（连交叉项都没有时，$f$ 本身已是 $n-1$ 元二次型，直接归纳。）
\textbf{考场动作}：情形一不必固定配 $x_1$，优先选"有平方项且系数最简单"的那个变量；
情形二先作上面的和差替换，把交叉项拆成平方项再回到情形一。

【反哺点】：服务考纲〔二次型〕"会用配方法化二次型为标准形"。配方是\textbf{唯一不必求特征值、不必正交化}的方法，纯代数操作，速度最快。学生的唯一卡点就是"二次型不含平方项"——记住 $x_1=z_1+z_2,\ x_2=z_1-z_2$ 这一步（其行列式为 $-2\ne0$，必非退化），任何二次型都能拉回情形一。
\end{fdeep}

% ---------- 初等变换法 ----------
% 来源：樊 8.2.3(2)（印刷页 512）；北大 五.2 定理 2 的矩阵证明（印刷页 145–147）
\begin{fdeep}{初等变换法：一行一列成对变换，同时记录变换矩阵}
除配方法与正交变换法之外的第三条路：对 $n$ 阶对称矩阵 $A$ 施行一系列\textbf{成对}的初等变换
——作一次列变换就立刻作一次同样的行变换——这正好是合同变换 $P^{\mathrm T}AP$。
把它写成上下两块的分块矩阵，全程只需盯住上块：
\fml{\begin{pmatrix}A\\E\end{pmatrix}\ \longrightarrow\ \begin{pmatrix}\Lambda\\C\end{pmatrix}}
（其中 $\Lambda$ 为对角阵，$C$ 为可逆阵）
操作规则：上块 $A$ 每作一次\textbf{列}变换，下块 $E$ 就跟着作一次\textbf{同样的列}变换；
上块再补作一次对应的\textbf{行}变换（下块不动）。做完后 $C^{\mathrm T}AC=\Lambda$，
令 $x=Cy$ 即得标准形 $f=y^{\mathrm T}\Lambda y=\lambda_1y_1^{2}+\dots+\lambda_ny_n^{2}$。
\textbf{与配方法的差别}：配方法只交出替换、不管矩阵；初等变换法把\textbf{可逆矩阵 $C$ 一并算出来}，
所以题目一旦要求"写出所作的非退化线性替换"，它比配方更稳、更不易出错。

【反哺点】：服务考纲〔二次型〕"会用配方法化二次型为标准形""掌握用正交变换化二次型为标准形"。当题干要求\textbf{同时给出标准形与所用替换 $x=Cy$}（$C$ 只要求可逆、不必正交）时，初等变换法是标准解法：下块忠实记录了每一步，$C$ 是"白送"的。唯一要守住的一条纪律——\textbf{下块永远只作列变换}，写成行变换就前功尽弃。
\end{fdeep}

% ---------- 顺序主子式法（Jacobi） ----------
% 来源：樊 例 8.34（印刷页 521–522）
\begin{fdeep}{顺序主子式法：不求特征值也能写出标准形}
设 $n$ 阶实对称矩阵 $A$ 的各阶顺序主子式 $\Delta_1,\Delta_2,\dots,\Delta_n$ \textbf{全不为零}
（约定 $\Delta_0=1$）。则存在\textbf{非退化}线性替换 $X=CY$ 使
\fml{f=\frac{\Delta_1}{\Delta_0}y_1^{2}+\frac{\Delta_2}{\Delta_1}y_2^{2}
+\dots+\frac{\Delta_n}{\Delta_{n-1}}y_n^{2}}
即第 $k$ 个系数恰为 $\lambda_k=\Delta_k/\Delta_{k-1}$，并且还可要求 $C$ 是单位上三角矩阵。
证明思路（对阶数归纳）：把 $A$ 分块为 $\begin{pmatrix}A_{n-1}&\beta\\ \beta^{\mathrm T}&a_{nn}\end{pmatrix}$，
由归纳假设有 $S$ 使 $S^{\mathrm T}A_{n-1}S=\operatorname{diag}(\lambda_1,\dots,\lambda_{n-1})$；
再用一个单位上三角的合同变换消去右上角 $\beta$，即得对角阵
$\operatorname{diag}(\lambda_1,\dots,\lambda_{n-1},\lambda_n)$，两边取行列式就有
$\lambda_n=\Delta_n/\Delta_{n-1}$。
\textbf{使用前提}：所有 $\Delta_k\ne0$；若某个 $\Delta_k=0$，本法失效，须先作一次合同置换再试。

【反哺点】：服务考纲〔二次型〕"掌握用正交变换化二次型为标准形，会用配方法化二次型为标准形"。这是\textbf{第三种、也是最省事}的求标准形办法：只算 $n$ 个顺序主子式，不用求特征值、不用求特征向量、不用配方；而 $\Delta_k/\Delta_{k-1}$ 的符号直接给出正负惯性指数（看相邻两个 $\Delta$ 是否同号），判正定、判合同一步到位。
\end{fdeep}

% ---------- A 与 B 合同的判定 ----------
% 来源：樊 8.2.2(3)（印刷页 511）；北大 五.2 定理 2、五.1（印刷页 140、145）
\begin{fdeep}{$A$ 与 $B$ 合同的判定：比正负惯性指数}
\textbf{核心判据}：$n$ 阶实对称矩阵 $A$ 与 $B$ \textbf{合同} $\iff$ 它们的正惯性指数相同
\textbf{且}秩相同；等价说法是"对应的二次型有相同的规范形"（两矩阵必合同于
$\operatorname{diag}(E_p,-E_{r-p},O)$，全部合同类由 $(p,\ r-p)$ 分类）。
考场三步：先比\textbf{秩}（秩不等直接不合同，因为合同必等价、等价保秩）；
秩相等再比\textbf{正惯性指数 $p$}（用顺序主子式法数正系数，或数正特征值个数）；
两项都相等则合同，否则不合同。
另有一条常考的唯一性结论：若 $A^{\mathrm T}=A,\ B^{\mathrm T}=B$ 且对一切 $x$ 有
\fml{x^{\mathrm T}Ax=x^{\mathrm T}Bx\ \Longrightarrow\ A=B}
（二次型与对称矩阵\textbf{一一对应}：$x_ix_j\ (i\ne j)$ 项系数的一半即 $a_{ij}$，$x_i^{2}$ 项系数即 $a_{ii}$。）

【反哺点】：服务考纲〔二次型〕"了解合同变换与合同矩阵""了解…惯性定理"。选择题"下列矩阵中与 $A$ 合同的是"，标准解法就是\textbf{算秩 + 算正惯性指数}这两个数，根本不必求出变换矩阵。再加一句"对称矩阵由二次型唯一决定"，可秒杀一类"已知 $x^{\mathrm T}Ax=x^{\mathrm T}Bx$，证明 $A=B$"的证明题。
\end{fdeep}

% ---------- 惯性指数 = 特征值的符号计数 ----------
% 来源：樊 8.2.2(4)（印刷页 511）
\begin{fdeep}{惯性指数就是特征值的符号计数}
对实对称矩阵 $A$（$r=\operatorname{rank}A$）：\textbf{正惯性指数 $p$ 等于正特征值的个数}，
\textbf{负惯性指数 $r-p$ 等于负特征值的个数}，秩等于非零特征值的个数。理由是正交变换下
\fml{Q^{\mathrm T}AQ=Q^{-1}AQ=\operatorname{diag}(\lambda_1,\dots,\lambda_n)}
它\textbf{既是合同又是相似}：合同保惯性指数、相似保特征值，两边的 $\pm1$ 计数必须对上。
由此得几条好用的推论：
（1）$A$ 正定 $\iff$ 全部特征值 $>0$；半正定 $\iff$ 全部特征值 $\ge0$（且至少有一个为 $0$）；
（2）符号差 $=2p-r=\sum_{i=1}^{n}\operatorname{sgn}\lambda_i$；
（3）正负惯性指数可"数特征值符号"，也可"看 $\Delta_k/\Delta_{k-1}$ 的符号"（即相邻两个顺序主子式
是否同号），\textbf{两条路算出的答案必须一致}。

【反哺点】：服务考纲〔二次型〕"掌握二次型及其矩阵表示""掌握用正交变换化二次型为标准形"，并与〔矩阵的特征值和特征向量〕打通。求正负惯性指数有三条路：数正特征值、看顺序主子式是否变号、配方法数正系数。\textbf{用两条不同的路算同一个 $p$}，是二次型部分最便宜的错误拦截手段（若用正交变换算出的标准形系数之和与 $\operatorname{tr}A$ 不符，也立刻暴露特征值算错）。
\end{fdeep}

% ---------- 二次型的秩 ----------
% 来源：樊 例 8.22（印刷页 514–516）；北大 五.3（印刷页 148）
\begin{fdeep}{二次型的秩：三种算法与 rank$(A^{\mathrm T}A)$}
二次型的\textbf{秩}定义为它矩阵 $A$ 的秩，也等于标准形（规范形）中\textbf{非零平方项的个数}
——合同变换保秩，而对角阵的秩就是非零对角元的个数。由此得三条实用结论：
（1）秩是合同不变量中最弱的一个：\textbf{秩不同，一定不等价、不合同、不相似}；
（2）若 $f=\sum_{i=1}^{s}\bigl(a_{i1}x_1+a_{i2}x_2+\dots+a_{in}x_n\bigr)^{2}$，
记系数矩阵为 $A_{s\times n}$，则 $f$ 的矩阵是 $A^{\mathrm T}A$，而
\fml{\operatorname{rank}(A^{\mathrm T}A)=\operatorname{rank}A}
故 $f$ 的秩 $=\operatorname{rank}A\le\min\{s,n\}$；
（3）特别地，$s<n$ 时 $f$ \textbf{不可能正定}（最多半正定）——正定要求有 $n$ 个平方项
且系数矩阵满秩。

【反哺点】：服务考纲〔二次型〕"了解二次型秩"。"$f$ 是若干个完全平方之和"是命题人最爱的给法：一眼写出矩阵 $A^{\mathrm T}A$，秩就等于系数矩阵的秩，比展开再求秩快得多；同时它解释了为什么"$n$ 元二次型写成少于 $n$ 个平方和"必然不正定——半正定与正定的分水岭就在这里。
\end{fdeep}

% ---------- 惯性指数的估计 ----------
% 来源：樊 例 8.36（印刷页 522–523）
\begin{fdeep}{惯性指数的估计：先把二次型写成 $p$ 个平方减 $q$ 个平方}
设 $f=l_1^{2}+\dots+l_p^{2}-l_{p+1}^{2}-\dots-l_{p+q}^{2}$，其中每个 $l_i$ 都是
$x_1,\dots,x_n$ 的\textbf{一次齐次式}。则 $f$ 的\textbf{正惯性指数 $\le p$，负惯性指数 $\le q$}。
证明思路（反证 + 齐次方程组，与惯性定理同一套路）：设 $f$ 经 $X=TY$ 化为标准形
$y_1^{2}+\dots+y_r^{2}-y_{r+1}^{2}-\dots$，若 $r>p$，则方程组
（前 $p$ 个取 $l_i=0$，后 $n-r$ 个取 $y_j=0$，共 $p+(n-r)<n$ 个方程）必有非零解 $X_0$。
此时左端 $f(X_0)=-l_{p+1}^{2}-\dots-l_{p+q}^{2}\le0$，右端却是 $y_1^{2}+\dots+y_r^{2}>0$，矛盾。
\textbf{用途}：数一数已给的平方项个数，就能立刻给出正惯性指数的\textbf{上界}。

【反哺点】：服务考纲〔二次型〕"理解正定二次型、正定矩阵的概念，并掌握其判别法"。见到 $f=\sum(\ \cdot\ )^{2}-\sum(\ \cdot\ )^{2}$ 这种结构（含参数题尤其常见），先用本结论卡住 $p$ 的上界，往往一步就排除一半选项；它也是判断"$f$ 是否可能正定"最快的反驳工具（平方项不足 $n$ 个即不正定）。
\end{fdeep}

% ---------- 四种定性的判别 ----------
% 来源：北大 五.4 定义 7 及其后（印刷页 155–156）
\begin{fdeep}{负定、半正定、半负定、不定的判别}
设 $f(x)=x^{\mathrm T}Ax$（$A$ 实对称，$x\ne0$）：$f>0$ 称\textbf{正定}；$f<0$ 称\textbf{负定}；
$f\ge0$ 称\textbf{半正定}；$f\le0$ 称\textbf{半负定}；既有 $f(x_1)>0$ 又有 $f(x_2)<0$ 称\textbf{不定}。
\textbf{负定没有独立的判别法}——加个负号化成正定即可：$f$ 负定 $\iff$ $-f$ 正定 $\iff$ $-A$ 正定，
于是它的顺序主子式满足
\fml{(-1)^{k}\Delta_k>0\qquad(k=1,2,\dots,n)}
即\textbf{奇数阶顺序主子式为负、偶数阶为正}。
\textbf{不定}的判定最省事：只要能举出一个 $x$ 使 $f(x)>0$、另一个使 $f(x)<0$ 即可；
或看特征值有正有负、正负惯性指数都不为 $0$。此外还有一条好用的等价说法：
$A$ 半正定或半负定 $\iff$ 凡满足 $\alpha^{\mathrm T}A\alpha=0$ 的 $\alpha$ 都有 $A\alpha=0$
（必要条件显然；充分性用反证：若正负惯性指数都 $\ge1$，取规范形下 $y_{p+1}=1$、其余为 $0$，
得 $\alpha^{\mathrm T}A\alpha=0$ 而 $A\alpha\ne0$）。

【反哺点】：服务考纲〔二次型〕"理解正定二次型、正定矩阵的概念，并掌握其判别法"。考纲字面只写"正定"，但选择题常考负定与不定。两条口诀：\textbf{负定 = 加负号变正定，故判据为 $(-1)^{k}\Delta_k>0$}；\textbf{判不定只需举两个反例向量}。最后那条"$\alpha^{\mathrm T}A\alpha=0\Longrightarrow A\alpha=0$"是判半正定/半负定的巧法。
\end{fdeep}

% ---------- 半正定的判别（含反例） ----------
% 来源：北大 五.4 定理 8（印刷页 156）；樊 8.3.2、例 8.82（印刷页 526、545）
\begin{fdeep}{半正定的判别：为什么"顺序主子式非负"不够}
$n$ 元实二次型 $f=x^{\mathrm T}Ax$（$A$ 实对称）半正定，有以下\textbf{等价条件}：
（1）$f$ 半正定；（2）正惯性指数 $p$ 等于秩 $r$（即负惯性指数为 $0$）；
（3）存在可逆 $C$ 使 $C^{\mathrm T}AC=\operatorname{diag}(d_1,\dots,d_n)$，其中 $d_i\ge0$；
（4）存在实矩阵 $C$ 使 $A=C^{\mathrm T}C$；
（5）$A$ 的\textbf{所有主子式}（行指标与列指标相同）都 $\ge0$。
\textbf{关键差别}：正定只需查\textbf{顺序}主子式，半正定必须查\textbf{所有}主子式。
\textbf{反例}：$f(x_1,x_2)=-x_2^{2}$，其矩阵为
$\begin{pmatrix}0&0\\0&-1\end{pmatrix}$，顺序主子式为 $0,\ 0$ 全部 $\ge0$，
但 $f(0,1)=-1<0$，\textbf{不是半正定}——漏掉的是一阶主子式 $a_{22}=-1$。
（证明思路：必要性取 $x$ 只在前 $k$ 个分量上变动；充分性反证，若正负惯性指数都 $\ge1$，
取 $C^{\mathrm T}AC=\operatorname{diag}(E_p,-E_q,O)$ 下 $y_{p+1}=1$、其余为 $0$ 的那个向量。）

【反哺点】：服务考纲〔二次型〕"掌握正定矩阵的判别法"（半正定是其自然延伸）。这是二次型部分\textbf{最经典的陷阱}：题干给"顺序主子式都 $\ge0$"，一定要警惕它只对正定充分。一句话记牢：\textbf{正定看顺序主子式，半正定看全部主子式}；判半正定还有个更快的写法——把它写成 $C^{\mathrm T}C$ 或若干平方和。
\end{fdeep}

% ---------- 分解为两个一次型之积 ----------
% 来源：樊 例 8.33（印刷页 521）
\begin{fdeep}{二次型能分解为两个一次型之积的充要条件}
非零实二次型 $f(x_1,\dots,x_n)$ 可以写成
\fml{f=(u_1x_1+\dots+u_nx_n)(v_1x_1+\dots+v_nx_n)}
的\textbf{充要条件}是：$f$ 的\textbf{秩为 $1$}，或者\textbf{秩为 $2$ 且符号差为 $0$}。
换成惯性指数的说法：正负惯性指数只能是 $(1,0)$、$(0,1)$（秩 $1$，即 $f$ 是正/负定的一次型平方）
或 $(1,1)$（秩 $2$ 且符号差 $0$）。证明思路：若 $u,v$ 线性相关，不妨 $u=kv$，
作替换 $y_1=v_1x_1+\dots+v_nx_n,\ y_i=x_i\ (i\ge2)$ 得 $f=ky_1^{2}$，秩为 $1$；
若 $u,v$ 线性无关，令 $y_1=\sum u_ix_i,\ y_2=\sum v_ix_i$ 得 $f=y_1y_2$，
再令 $y_1=z_1+z_2,\ y_2=z_1-z_2$ 得 $f=z_1^{2}-z_2^{2}$，正是秩 $2$、符号差 $0$。
反之把 $y$ 用 $C^{-1}$ 反解回 $x$ 即可构造出两个一次型。

【反哺点】：服务考纲〔二次型〕"了解二次型秩…及惯性定理"。"$f$ 能否分解为两个一次因式之积"是反复出现的判断题，硬展开系数矩阵很慢；\textbf{只要求出秩与符号差，答案立刻出来}：秩 $1$ 行；秩 $2$ 看符号差是否为 $0$；秩 $\ge3$ 一定不行。这是惯性指数最"解题化"的一个用处。
\end{fdeep}

% ---------- 同时合同对角化 ----------
% 来源：樊 例 8.67、例 8.35、例 8.66（印刷页 521、537–538）
\begin{fdeep}{同时合同对角化：$A$ 正定时可以把 $B$ 一起化为对角阵}
设 $A$ 是 $n$ 阶\textbf{正定}矩阵，$B$ 是 $n$ 阶\textbf{实对称}矩阵，则存在可逆矩阵 $P$ 使
\fml{P^{\mathrm T}AP=E,\qquad P^{\mathrm T}BP=\operatorname{diag}(\mu_1,\dots,\mu_n)}
其中 $\mu_1,\dots,\mu_n$ 恰是方程 $|\lambda A-B|=0$ 的 $n$ 个根，且全为实数。
证法：$A$ 正定 $\Longrightarrow$ 存在可逆 $C$ 使 $C^{\mathrm T}AC=E$；此时 $C^{\mathrm T}BC$ 仍实对称，
故存在\textbf{正交}矩阵 $Q$ 使 $Q^{\mathrm T}(C^{\mathrm T}BC)Q=\operatorname{diag}(\mu_i)$，
取 $P=CQ$（正交变换不会破坏已化好的 $E$）。两个常用推论：
（1）$A$ 正定、$B$ 实对称时，存在可逆变换 $x=Py$ 把 $x^{\mathrm T}Ax$ 化为\textbf{规范形}、
同时把 $x^{\mathrm T}Bx$ 化为\textbf{标准形}；
（2）若 $|A-\lambda B|=0$ 有 $n$ 个\textbf{互异}实根，则存在非退化 $X=CY$ 同时把两个二次型化为标准形。

【反哺点】：服务考纲〔二次型〕"掌握用正交变换化二次型为标准形的方法"。"两型同时化简"是二次型大题的常见收尾要求。固定套路是\textbf{先把正定的那个 $A$ 化为 $E$}（用配方或初等变换都行），再对剩下的实对称阵作正交对角化；顺序反了就做不出来，$\mu_i$ 也一定要交代清楚是 $|\lambda A-B|=0$ 的根。
\end{fdeep}

% ---------- 正定矩阵的元素与行列式结论 ----------
% 来源：樊 例 8.63、例 8.75、例 8.25、例 8.50、例 8.51（印刷页 528、541–542、516、531）
\begin{fdeep}{正定矩阵的元素、行列式与合同结论}
$A$ 正定（实对称）时，以下都是可以\textbf{直接引用}的结论：
（1）$|A|>0$，所有对角元 $a_{ii}>0$；更强的有
\fml{|A|\le a_{11}a_{22}\cdots a_{nn},\qquad |A|\le a_{nn}\Delta_{n-1}}
（$\Delta_{n-1}$ 为 $n-1$ 阶顺序主子式）；
（2）$A$ 中\textbf{绝对值最大的元素必在主对角线上}：由二阶主子式 $a_{ii}a_{jj}-a_{ij}^{2}>0$
得 $|a_{ij}|<\sqrt{a_{ii}a_{jj}}\le a_{kk}$，其中 $a_{kk}=\max\{a_{11},\dots,a_{nn}\}$；
（3）$A$ 与 $A^{-1}$ \textbf{合同}（因 $A^{-1}=(A^{-1})^{\mathrm T}AA^{-1}$），
所以 $x^{\mathrm T}Ax$ 与 $x^{\mathrm T}A^{-1}x$ 有相同的正负惯性指数；
（4）对 $m\times n$ 实矩阵 $B$（$m\ge n$），$B^{\mathrm T}AB$ 正定 $\iff$ $\operatorname{rank}B=n$；
（5）$A$ 可逆 $\iff$ 存在矩阵 $B$ 使 $AB+B^{\mathrm T}A$ 正定。

【反哺点】：服务考纲〔二次型〕"理解正定二次型、正定矩阵的概念，并掌握其判别法"，并与〔矩阵〕的逆与秩打通。(2) 把"找最大元"这类选择题一步解决；(3)(4) 是"抽象二次型判正定"的固定套路——见 $B^{\mathrm T}AB$ 就想 $\operatorname{rank}B=n$；(1) 是行列式估计题的入口。
\end{fdeep}

% ---------- 半正定的行列式单调性与次可加性 ----------
% 来源：樊 例 8.39、例 8.72、例 8.69（印刷页 524、540–539）
\begin{fdeep}{半正定的行列式单调性与惯性指数的次可加性}
（1）\textbf{惯性指数次可加}：$A,B$ 为 $n$ 阶实对称矩阵时
\fml{p(A+B)\le p(A)+p(B)}
（$p$ 为正惯性指数。半正定情形显然：$p(A+B)=\operatorname{rank}(A+B)\le\operatorname{rank}A+\operatorname{rank}B$；
一般情形把 $A$ 拆成半正定阵与半负定阵之和、$B$ 同样拆分，再用半正定情形相加。）
（2）\textbf{行列式单调性}：若 $A,B,A-B$ 都半正定（记 $A\ge B\ge0$），则 $\det A\ge\det B\ge0$；
若 $B$ 正定而 $A-B$ 半正定，则 $|A|\ge|B|$，且 $|\lambda A-B|=0$ 的根全满足 $\lambda\ge1$。
（3）$A,B$ 都正定 $\Longrightarrow |A+B|\ge|A|+|B|$。

【反哺点】：服务考纲〔二次型〕"掌握正定矩阵的判别法"。这三条是"用正定性证不等式"的标准工具：(1) 用来估计含参数二次型的正惯性指数上界；(2)(3) 用来证 $\det$ 型不等式（如"$A$ 正定、$A-B$ 半正定，证 $|A|\ge|B|$"），比逐个估计特征值快得多，也是"比较两个正定矩阵大小"的通用语言。
\end{fdeep}

% ---------- 正定矩阵的 Cauchy 型不等式 ----------
% 来源：樊 例 8.76、例 8.77（印刷页 542–543）
\begin{fdeep}{正定矩阵的 Cauchy 型不等式}
设 $A$ 为 $n$ 阶正定矩阵，$\alpha,\beta,x,y\in\mathbb R^{n}$，则
\fml{(\alpha^{\mathrm T}A\beta)^{2}\le(\alpha^{\mathrm T}A\alpha)(\beta^{\mathrm T}A\beta)}
等号成立当且仅当 $\alpha,\beta$ \textbf{线性相关}。证明只需把 $A$ 写成 $A=Q^{\mathrm T}Q$（$Q$ 可逆），
则 $\alpha^{\mathrm T}A\beta=(Q\alpha)^{\mathrm T}(Q\beta)$，对两个向量 $Q\alpha,Q\beta$ 用通常的
Cauchy 不等式即可。把 $A$ 换成 $A^{-1}$ 又得一组对偶形式：
\fml{(x^{\mathrm T}y)^{2}\le(x^{\mathrm T}Ax)(y^{\mathrm T}A^{-1}y),\qquad
x^{\mathrm T}y\le\frac12\bigl(x^{\mathrm T}Ax+y^{\mathrm T}A^{-1}y\bigr)}
后者的等号成立当且仅当 $y=Ax$（由 $(x-A^{-1}y)^{\mathrm T}A(x-A^{-1}y)\ge0$ 展开即得）。

【反哺点】：服务考纲〔二次型〕"掌握正定矩阵的判别法"，并与〔向量〕的 Cauchy 不等式呼应。"已知 $A$ 正定，证明某个二次型不等式"的题，几乎都是\textbf{把 $A$ 写成 $Q^{\mathrm T}Q$ 后套Cauchy 不等式}；这两组不等式及其取等条件要正反两用——既用来证不等式，也用来由等号成立反推出"向量线性相关"或"$y=Ax$"。
\end{fdeep}

% ---------- 两类高频特殊二次型 ----------
% 来源：樊 例 8.19、例 8.20（印刷页 513–514）
\begin{fdeep}{两类高频特殊二次型的秩与符号差}
（1）$f=ayz+bzx+cxy$（$a,b,c$ 为实数）：
$a,b,c$ 全为零时，秩 $=0$、符号差 $=0$；$abc=0$ 但不全为零时，秩 $=2$、符号差 $=0$；
$abc>0$ 时，秩 $=3$、符号差 $=-1$；$abc<0$ 时，秩 $=3$、符号差 $=1$。
（推导：$c\ne0$ 时 $A$ 合同于 $\operatorname{diag}\bigl(c,-c,-\frac{ab}{c}\bigr)$，再按 $c$ 与 $ab$ 的符号分类。）
（2）$f=n\sum_{i=1}^{n}x_i^{2}-\Bigl(\sum_{i=1}^{n}x_i\Bigr)^{2}$，它还可以写成离差平方和
\fml{f=\sum_{1\le i<j\le n}(x_i-x_j)^{2}}
所以它是\textbf{半正定}二次型，秩 $=$ 正惯性指数 $=n-1$，负惯性指数为 $0$，
规范形是 $y_1^{2}+\dots+y_{n-1}^{2}$。（它的矩阵 $A=nE-\alpha\alpha^{\mathrm T}$，
$\alpha=(1,1,\dots,1)^{\mathrm T}$，特征值为 $n$（$n-1$ 重）与 $0$。）

【反哺点】：服务考纲〔二次型〕"了解二次型秩…及定义与惯性定理"。这两类结构在填空、选择里反复出现："$axy+byz+czx$ 的秩与符号差"按 $abc$ 的符号\textbf{直接背出四种答案}；"$n\sum x_i^{2}-(\sum x_i)^{2}$"要能一眼看出它是离差平方和，从而半正定、秩 $n-1$、负惯性指数 $0$，完全不必展开求特征值。
\end{fdeep}


% ===========================================================================
%  【主控裁决 · 第 9 章 deep_ch09 取舍】（依据 AGENTS.md §2 决定二）
%
%  第 9 讲是全书最长的一讲（9 讲 18 页），故深化料可留较厚；但仍按「能否反哺解题」筛一遍。
%
%  【保留 15 块】
%   ✓ 块1　{惯性定理的证明思路：为什么正负惯性指数是不变量
%   ✓ 块2　{规范形唯一、标准形不唯一：复域与实域两套规范形
%   ✓ 块3　{配方法的两种情形与操作步骤
%   ✓ 块4　{初等变换法：一行一列成对变换，同时记录变换矩阵
%   ✓ 块5　{顺序主子式法：不求特征值也能写出标准形
%   ✓ 块6　{$A$ 与 $B$ 合同的判定：比正负惯性指数
%   ✓ 块7　{惯性指数就是特征值的符号计数
%   ✓ 块8　{二次型的秩：三种算法与 rank$(A^{\mathrm T}A)$
%   ✓ 块9　{惯性指数的估计：先把二次型写成 $p$ 个平方减 $q$ 个平方
%   ✓ 块10　{负定、半正定、半负定、不定的判别
%   ✓ 块11　{半正定的判别：为什么"顺序主子式非负"不够
%   ✓ 块12　{二次型能分解为两个一次型之积的充要条件
%   ✓ 块13　{同时合同对角化：$A$ 正定时可以把 $B$ 一起化为对角阵
%   ✓ 块14　{正定矩阵的元素、行列式与合同结论
%   ✓ 块17　{两类高频特殊二次型的秩与符号差
%
%  【删除 2 块及理由】
%   ✗ 块15　{半正定的行列式单调性与惯性指数的次可加性
%   ✗ 块16　{正定矩阵的 Cauchy 型不等式
%
%   · 块15「半正定的行列式单调性与次可加性」——惯性指数次可加性 $p(A+B)\le p(A)+p(B)$
%       与 $|A+B|\ge|A|+|B|$ 属纯理论拓展，考纲无接口，反哺度低。
%   · 块16「正定矩阵的 Cauchy 型不等式」——$(\alpha^{\mathrm T}A\beta)^2\le(\alpha^{\mathrm T}A\alpha)(\beta^{\mathrm T}A\beta)$
%       属推广型不等式，考研不考。
%
%  【保留重点（对齐考纲〔二次型〕3 条）】
%   惯性定理证明思路、规范形唯一性、**配方法两种情形**、**初等变换法**、**Jacobi 顺序主子式法**、
%   **合同判定**、惯性指数↔特征值符号、二次型的秩、惯性指数估计、
%   **负定/半正定/半负定/不定判别**、**半正定需查全部主子式的反例**、
%   分解为两一次型之积、**同时合同对角化**、正定矩阵元素与行列式结论、两类高频特殊二次型。
% ===========================================================================
